\section{A polarized representation on the completion at a two-block point}

% BEGIN U71R1
\paragraph{U71R1. The finite algebra over the original coefficient ring.}
Take the data in $\mathrm{Pa.S6}$. Write $H=G.G_0$ for the
index-two group of the source frame, $D=u.D$, $k=B/\mathfrak m_B$,
$A=x.K[[t]]$, $F=\operatorname{Frac}A$,
$C=B_h$ and $\mathfrak m=\mathfrak m_C$. The given map
$f=P.f:B\longrightarrow A$ sends $P.s$ to $t$.
The given $r_F:C\longrightarrow F$ agrees on $B$ with $f$
followed by inclusion and has kernel $\mathfrak m$.
Put $a=x.a$, so $1\leq a\leq n$.
The second block has size $n-a$ and may be empty.
No Selmer bound is a hypothesis of the construction.

Let $Z$ consist of the ordinary products, the elements $[g]-1$
for $g\in r.I_p(w)$, the two auxiliary families $q_1,q_2$,
and their polarized adjoints. The ideal of the point is
\[
 I=P.\mathrm{ker}=
 \bigcap_{j,V}
 \bigl(\mathrm{CH}(D),Z,\mathfrak m_B^j,[v]-1\ (v\in V)\bigr)
 \ \subseteq B[H],
\]
where $V$ ranges over open normal subgroups, the ideals are
two-sided and $\mathrm{CH}(D)$ uses all coefficients of the generic
Cayley--Hamilton identities. Write $E_B=B[H]/I$.

The residue field of $B$ equals that of $O$ and is finite.
The residual determinant is the product of the distinct residue
characters of $r$. SCI2--3 thus apply in their multiplicity-free
form. The equations $u.\mathrm{Ord},u.I,P.q_1,P.q_2$ give
$\operatorname{tr}_D([h]z)=0$ for every $h\in H$ and every
unstarred relation $z$ in $Z$. The adjoint is the $B$-linear
anti-involution defined on group elements by
\[
 [g]^*=\varepsilon(g)^{1-n}[\mathrm{conj}(g)^{-1}].
\]
The characteristic coefficients of a group element and of its
adjoint agree by $u.P$. The integral Amitsur expansion (Chenevier, Lemma~1.12) for a
generic group-algebra sum extends this equality to the full
polynomial law: reversing a primitive cyclic word permutes the
cyclic words in the expansion, and the scalar factors of its
letters multiply. The coefficient monomial and the exponent of each
characteristic coefficient are preserved in this permutation.
In particular the trace, the full kernel, and the generic
Cayley--Hamilton ideal are invariant under the adjoint.
SCI3 puts all of $Z$ in the full kernel. PQL7--8 then put $I$
in that kernel. The involution preserves the displayed intersection:
in each open-normal-subgroup test one may shrink the subgroup so that
its conjugate lies in $V$ and $\varepsilon^{1-n}$ is $1$ modulo
$\mathfrak m_B^j$ there. Thus $E_B$ has the induced involution and
a Cayley--Hamilton determinant with cyclic, adjoint-invariant trace.
All Cayley--Hamilton identities continue to hold after scalar
extension by the generic coefficient construction.

Here is a finite generation argument which also controls the topology.
Apply SCI2 first to the \emph{algebraic} quotient
$E'=B[H]/(\mathrm{CH}(D))$. Its complete orthogonal idempotents
$p_1,\ldots,p_n$ satisfy $p_iE'p_i=Bp_i$. These idempotents
concern the original residue characters over $k$.
The finite topological generating set in $P.\mathrm{fg}$ gives
one for $H$: the dense abstract subgroup in the ambient group
meets the open index-two kernel densely, and Schreier's coset
generators generate that intersection. Take a finite symmetric
list $S$ of generators of this dense subgroup of $H$.
Expand a word in $S$ by inserting $\sum p_i=1$ between its
letters and at both ends. A path which visits the same vertex
twice contains a closed subpath in $p_iE'p_i$. Replace it by
its scalar times $p_i$ and shorten the path. Consequently the
span of the paths of length at most $n-1$ in these finitely many
letters and idempotents contains every such word in $E'$.

All those finitely many paths have representatives in $B[H]$.
The quotients of $B[H]$ by the individual ideals in the
intersection for $I$ are finite: they are quotients of
$(B/\mathfrak m_B^j)[H/V]$. The paths span each quotient,
because the chosen dense subgroup maps onto $H/V$. They also
span the inverse limit. Indeed, $B$ is compact as the inverse
limit of its finite adic quotients. For a compatible element of
the limit, the tuples of coefficients on our fixed list which
give that element at any one level form nonempty closed subsets
of a finite product of copies of $B$. Every finite set of levels
has a common refinement. Compactness gives one tuple working
at all levels. The kernel of $B[H]$ in the limit is exactly $I$,
and the path representatives map onto the limit.
This identifies the algebraic $E_B$ with the profinite limit
and proves it is finite over $B$.

The adic topology of this finite module is the inverse-limit
topology. Indeed the map from the finite product of copies of
$B$ on the paths is continuous onto the Hausdorff limit, and
compactness makes its topology the quotient topology.
Its images of $\mathfrak m_B^j$ times that free module are
$\mathfrak m_B^j E_B$. The map $H\to E_B$ is continuous by
its finite-quotient construction. For each $j$ choose an open
normal $V$ such that $[v]-1\in\mathfrak m_B^jE_B$ for $v\in V$.
Multiplication by $[g]$ gives the same bound on $[gv]-[g]$
simultaneously for all $g$.
% END U71R1

% BEGIN U71R2
\paragraph{U71R2. The Laurent specialization and self-adjoint coordinate projections.}
Put $E=C\otimes_B E_B$. It is finite over the complete
Noetherian local ring $C$ and carries the base-changed
Cayley--Hamilton determinant $\Delta$ and involution.
The given point matrices define an algebra map
$r:E\to M_n(F)$. To check that they factor, the ordered ordinary
product kills the whole local flag: the factor indexed by $i$
maps step $i$ into step $i-1$, starting with the rightmost
factor. At $r.I_p(w)$ the matrices are the identity by $P.I_p$
and $x.\mathrm{ine}$. The first auxiliary product vanishes by
$P.\mathrm{eig}$ and ordinary Cayley--Hamilton; the second
vanishes because $h.p_0(\tau_j)=1$ and $f(y_{ji})=0$.
The adjoints of these relations vanish by the polarized matrix
identity at $x$, since $h$ respects conjugation and cyclotomic
characters. The generic Cayley--Hamilton coefficients vanish for
the matrix determinant $D_f$ supplied by $P.\mathrm{maps}$.
Finally the matrix map into $M_n(A)$ kills the intersection
defining $I$. Indeed modulo $t^j$ the continuous group map
factors through $H/V$ for an open normal $V$ and
$f(\mathfrak m_B^j)\subseteq(t^j)$. Here $f$ is local since
the composite $B\to A\to x.K$ is surjective and has the unique
maximal ideal as kernel. Separatedness of $A$ gives the assertion.
Tensoring with $C$ gives the map $r$, with
\[
 r(y^*)=r(y)^t,\qquad
 \overline{\chi_\Delta(y)}=\det(X-r(y)).
\]
The second identity follows for arbitrary $C$-linear sums of
group elements from equality of the full laws after pullback
to $F$, by $P.\mathrm{maps}$.

In fact $f$ is onto: lift each coefficient of a power series
from $x.K$ to $O$ and sum its multiples of $s^j$ in $B$.
Here $s\in\mathfrak m_B$ because its image $t$ is a nonunit.
The image of every nonzero denominator in $A$ has a lift to
a unit of $C$ since $\ker r_F=\mathfrak m$.
Thus $r_F$ is surjective. Its residue field is $F$, of odd
characteristic, so $2$ is a unit of $C$.
The image of $r$ is
\[
 M_a(F)\times M_{n-a}(F)
\]
in the prescribed ordered coordinates (omit an empty factor).
Indeed $x.\mathrm{blk}$ gives one containment. The other follows
by clearing the denominators of matrix units and applying
$x.\mathrm{span}$ to them, then inverting $t$. Every group element
thus used lifts through $h.p_0$, since $h.p$ is onto and preserves
the index-two character.

Write $J=\ker r$. The ideal $J/\mathfrak m E$ is nilpotent.
For each $y\in J$, the displayed characteristic polynomial
reduces to $X^n$, so $y^n\in\mathfrak m E$. An ideal of nil
elements in a finite dimensional algebra over a field is contained
in its Jacobson radical: $1-zy$ is a unit for every $z$, by a
finite geometric series. The radical of the finite dimensional
algebra is nilpotent, by stabilization of its powers and Nakayama's
lemma. Thus $J^d\subseteq\mathfrak m E\subseteq J$ for some $d$,
and completeness of the finite $C$-module $E$ gives completeness
for the powers of $J$.

Lift the successive diagonal projections of $r(E)$ to orthogonal
self-adjoint idempotents $e_1,\ldots,e_n$. In the complement
corner of those already chosen, lift the next projection and
average the lift with its adjoint using $1/2$. For the resulting
self-adjoint $u$, the error $u^2-u$ lies in $J$. The polynomial
$3u^2-2u^3$ equals $u$ modulo that error and its new error is
divisible by $(u^2-u)^2$. Iteration converges to the required
idempotent in the same corner. Take the final complement itself.
Thus $\sum e_i=1$ and $r(e_i)$ is the specified $i$th diagonal
matrix unit. These projections need no particular relation to the
idempotents $p_i$ used over $B$.
% END U71R2

% BEGIN U71R3
\paragraph{U71R3. Rank one in the determinant and normalization within each block.}
For one $e=e_i$ consider
$P_e(X)=\Delta((1-e)+Xe)$. Multiplicativity and $P_e(1)=1$
imply that the coefficients of $P_e$ are orthogonal idempotents
with sum one: compare coefficients in $P_e(XY)=P_e(X)P_e(Y)$.
Since $C$ is local this polynomial is one monomial. Its reduction
is $X$ by U71R2 and equality of the specialized determinants.
Thus $P_e=X$. The corresponding polynomial for $1-e$ is
$X^{n-1}$, since the product of these two polynomials is
$\Delta(X\cdot1)=X^n$.

Define $\delta_e(y)=\Delta((1-e)+y)$ on $eEe$, with the same
definition on every scalar extension. It is multiplicative,
homogeneous of degree one by the equation for $P_e$, and
$\delta_e(e)=1$. A homogeneous polynomial law of degree one
is linear: on a finite tuple expand in independent scalar
variables and use the degree-one condition with an extra
variable. Only the individual linear terms survive. For $y\in
eEe$, apply Cayley--Hamilton to $y+u(1-e)$ over $C[u]$.
Its characteristic polynomial is
$(X-u)^{n-1}(X-\delta_e(y))$, by factoring into its two
complementary corners with identity on the other corner in
each factor. In the $e$ corner this gives
\[
 (y-ue)^{n-1}(y-\delta_e(y)e)=0.
\]
The leading coefficient in $u$ gives $y=\delta_e(y)e$.
Consequently $eEe=Ce$ freely on $e$; the scalar equals the
$\Delta$-trace, as seen from $\Delta(1+Ty)=1+T\delta_e(y)$.
The split corner descriptions persist after every scalar extension.

Choose a base index $o$ in each nonempty block. For an index
$i$ in its block choose $v_i\in e_oEe_i$ lifting the matrix unit
$E_{oi}$, using the surjection in U71R2 and the corner projections.
Then $v_i v_i^*=d_i e_o$ and $v_i^*v_i=d_i e_i$ with the same
scalar, by cyclic trace. Reduction gives $d_i\equiv1\bmod\mathfrak m$.
The root $1$ modulo $\mathfrak m$ of $X^2-d_i$ is simple.
Completeness gives a unique square root there by Hensel lifting.
Rescale $v_i$ by its inverse, so both products are the indicated
idempotents. Take $v_o=e_o$. The reduction of every $v_i$ is
still $E_{oi}$. Hence the coordinate of $y_{ij}=e_iye_j$ in
the corresponding base corner is
\[
                    v_i\,y_{ij}\,v_j^*.
\]
One recovers $y_{ij}$ by multiplication by $v_i^*$ and $v_j$.
Coordinates in a common block are scalars in $C$.
% END U71R3

% BEGIN U71R4
\paragraph{U71R4. A commutative coefficient algebra for both off-diagonal corners.}
Suppose first there are two blocks with base indices $o,p$.
Put $N=e_oEe_p$. Its adjoint identifies it $C$-linearly with
$e_pEe_o$. For $u,v\in N$ define $\beta(u,v)$ by
$uv^*=\beta(u,v)e_o$. This pairing is $C$-bilinear. Applying
the adjoint shows that it is symmetric. By cyclic trace,
$v^*u=\beta(u,v)e_p$. All its values belong to $\mathfrak m$:
$r$ kills $N$ because its two blocks are disjoint, and reducing
the defining multiplication takes the scalar to $r_F(\beta(u,v))$.
Associativity in $E$ gives, for $u,v,w\in N$,
\[
 \beta(u,v)w=(uv^*)w=u(v^*w)=\beta(v,w)u.
\]
Define $A_U=C\oplus N$ with unit $(1,0)$ and multiplication
\[
 (c,u)(d,v)=(cd+\beta(u,v),\,cv+du).
\]
This is a commutative associative ring. By bilinearity it suffices
for associativity to test three pure factors, and the only
remaining identity for three factors in $N$ was just proved.
It is a finite $C$-algebra, since $N$ is a direct summand of
the finite module $E$. The scalar inclusion is split and
$N^2\subseteq\mathfrak m$ inside this scalar copy.
With a single block set $A_U=C$ and $N=0$ instead.

The coordinates of U71R3 define a map
$\phi:E\to M_n(A_U)$. For an entry in the first-to-second
corner use its value in $N$; for the opposite corner use its
adjoint in $N$. Products through an intermediate index in one
block use scalar multiplication. Products which cross blocks
twice use $\beta$, in either order by the two scalar corner
identities above. Since $v_j^*v_j=e_j$, summing over $j$
then proves the matrix multiplication rule. The inverse on
entries from U71R3 gives a $C$-module isomorphism onto
\[
 \left\{(m_{ij}):\
   m_{ij}\in C\text{ within each block},\quad
   m_{ij}\in N\text{ across blocks}\right\}.
\]
Every set of allowed entries is attained at once by adding the
reconstructed corner elements. Also $\phi(1)=1$ and
$\phi(y^*)=\phi(y)^t$, by construction. Specialization of a
same-block scalar entry under $r_F$ gives the original entry of
$r(y)$, since the $v_i$ reduced to the prescribed matrix units.
Every cross-block entry of $r(y)$ is zero.
% END U71R4

% BEGIN U71R5
\paragraph{U71R5. Equality of the full determinant over arbitrary scalar algebras.}
We check that $\phi$ realizes $\Delta$ as a polynomial law. Let
$S$ be any commutative $A_U$-algebra, let $Y\in S\otimes_C E$,
and write $\phi_S$ for the induced map to $M_n(S)$.
The projections $e_i$ map to the diagonal matrix units. Diagonal
corners are still scalars, and
$\Delta(1-e_i+s e_i)=s$ after any scalar extension.
If $u=e_iue_j$ for $i\ne j$, the elementary operation $1+u$
has $\Delta$ equal to $1$: use the integral identity
$\Delta(1+e_i u)=\Delta(1+u e_i)$ from SCI3.
The matrix determinant of $\phi_S(1+u)$ is likewise one.
These identities apply over arbitrary rings with nilpotents.

Over $S[T]/T^{n+1}$ eliminate $1+TY$ using the coordinate
projections. The first diagonal entry is a scalar congruent
to $1$ modulo $T$, hence invertible. Left elementary operations
remove the other entries in its column, and right elementary
operations remove the other entries in its row. The remaining
corner matrix is still the identity modulo $T$. Repeating
gives a diagonal with invertible scalars $d_1,\ldots,d_n$.
The operations exist in the algebra itself: a column entry
in $e_i E_S e_j$ times the inverse scalar pivot is an admissible
off-diagonal element. In the matrix realization the same
operations perform the same elimination.
The diagonal is the product of the factors $1-e_i+d_i e_i$.
It follows that $\Delta(1+TY)$ and
$\det(1+T\phi_S(Y))$ both equal $\prod_i d_i$ in this truncated
polynomial ring. Both original polynomials have degree at most
$n$, so their top coefficients give
$\Delta_S(Y)=\det\phi_S(Y)$. Pulling this equality back to
$S[H]$ proves
\[
             D_{A_U}=\mathrm{GDet.ofRep}\bigl(g\mapsto\phi([g])\bigr).
\]
In particular this verification works in residue characteristics
at most $n$ and includes all mixed characteristic coefficients,
without recovering them from division in Newton identities.
% END U71R5

% BEGIN U71R6
\paragraph{U71R6. The prime test, the finite lattice, and all fields of $\mathrm{Pa.S6}$.}
Set $\rho_U(g)=\phi(1\otimes[g])$. This is a group representation
since the images of group elements multiply and are invertible.
Its $C$-linear span is the entire block module in U71R4:
$C[H]\to E$ is onto by the definition of $E$.
The identity $\phi(y^*)=\phi(y)^t$ and the group formula for the
adjoint give
\[
       \rho_U(g)^t\rho_U(\mathrm{conj}(g))=\varepsilon(g)^{1-n}1.
\]
If $\rho_U(g)_{ij}-b\in N$ for $b\in C$, a same-block entry
equals $b$ by $C\cap N=0$. Its reduction gives precisely the
$(i,j)$ entry of $x.\rho(h.p_0g)$ in $F$. For distinct blocks
the same hypothesis forces $b=0$, also giving that identity.
Every $z\in I$ maps to zero as a matrix by its image in $E_B$,
including the closure of all the defining relations.

For the thickness condition take any prime $Q\subset C$ and any
nonsemisimple witness $C\to\kappa$ with kernel $Q$ as in
$\mathrm{NonSS}$. If $s\in\operatorname{Ann}_C N$ were outside
$Q$, then $s$ would be invertible in $\kappa$ and
$\kappa\otimes_C N=0$. By the \emph{split module} description of
U71R4, the tensor of $E$ would be the product of the full matrix
algebras on the nonempty blocks over $\kappa$, with ordinary
matrix multiplication. On the other hand right exactness gives
\[
 \kappa\otimes_C E
 \simeq \kappa[H]/(\operatorname{image}(I)),
\]
where the image is taken through $B\to C\to\kappa$ and generates a
two-sided ideal. This is exactly the algebra in the witness, whereas
a product of full matrix algebras is semisimple. Thus
$\operatorname{Ann}_C N\subseteq Q$. This also covers $N=0$,
in which case such a witness cannot exist.

Finally take as the required $B$-module of matrices the image
of $E_B$ under $z\mapsto\phi(1\otimes z)$. It is finitely generated
over the \emph{original} $B$ by U71R1 and contains every $\rho_U(g)$.
For each $j$ the same open normal subgroup $V$ as in U71R1 ensures
$\rho_U(gv)-\rho_U(g)$ belongs to $\mathfrak m_B^j$ times this
module, for every $g\in H$ and $v\in V$.
The scalar tower is the composite $B\to C\to A_U$. U71R4 supplies
finiteness, scalar injectivity, the direct complement, its product
condition and the zero complement when $a=n$. U71R5 supplies the
required full determinant. Together these statements give the actual
$\mathrm{URep}$ and establish $\mathrm{Pa.s6}$ for every point and
completion in its original quantifiers.
% END U71R6
